%% kapa: exam=1 round=31 year=2020 period=2 form=A subject=accounting subject_ko=회계학 seq=39 number=79
%% source: Q-Net 감정평가사 기출문제 (cst00309) artlSeq=5207219 file=extracted/1008/5207219/A형/A형/2020년도 제31회 감정평가사 1차 시험_2교시_A형.pdf
%% conversion: kordoc 4.24.51 -> tools/gen1.py (auto); review=proofread-sonnet; flags=none
\begin{question}{79}
(주)감평은 최근 신제품을 개발하여 최초 10단위의 제품을 생산하는데 총 150시간의 노무시간을 소요하였으며, 직접노무시간당 ￦1,200의 직접노무원가가 발생하였다. (주)감평은 해당 신제품 생산의 경우, 90\%의 누적평균시간 학습곡선모형이 적용될 것으로 예상하고 있다. 최초 10단위 생산 후, 추가로 30단위를 생산하는데 발생할 것으로 예상되는 직접노무원가는?
\begin{choices}
  \item ￦180,000
  \item ￦259,200
  \item ￦324,000
  \item ￦403,200
  \item ￦583,200
\end{choices}
\end{question}
\answer{4}
